\documentclass[12pt]{article}

\input{../../_assets/preambulo.tex}


\begin{document}

    % 1. Foto de fondo
    % 2. Título
    % 3. Encabezado Izquierdo
    % 4. Color de fondo
    % 5. Coord x del titulo
    % 6. Coord y del titulo
    % 7. Fecha

    
    \input{../../_assets/portada}
    \portadaExamen{ffccA4.jpg}{Ecuaciones\\Diferenciales II\\Examen V}{Ecuaciones Diferenciales II. Examen V}{MidnightBlue}{-8}{28}{2026}{}

    \begin{description}
        \item[Asignatura] Ecuaciones Diferenciales II.
        \item[Curso Académico] 2022/2023.
        \item[Grado] Doble Grado en Ingeniería Informática y Matemáticas
        \item[Grupo] Único
        \item[Profesor] Rafael Ortega Ríos
        \item[Descripción] Parcial 1.
        \item[Fecha] 12 de Abril de 2023.
        % \item[Duración] ---.
    
    \end{description}
    \newpage


    % ------------------------------------
    
    \begin{ejercicio}
        Encuentra una función continua $x : [0, 2] \to \bb{R}$, $t \mapsto x(t)$, que cumpla
        \begin{equation*}
            x(t) = \begin{cases}
                3, & t \in [0, 1],\\
                \displaystyle 3 + 3 \int_{0}^{t-1} x(s)ds, & t \in [1, 2].
            \end{cases}
        \end{equation*}
        ¿Es única?
    \end{ejercicio}

    \begin{ejercicio}
        Encuentra una sucesión de funciones $\{f_n\}_{n \geq 0}$, con $f_n : [0, 1] \to \bb{R}^2$, que sea equicontinua pero no sea uniformemente acotada.
    \end{ejercicio}

    \begin{ejercicio}
        Demuestra que hay unicidad local para el problema de valores iniciales
        \[
            \dot{x} = \sqrt{|x|}, \quad x(2) = 4.
        \]
    \end{ejercicio}

    \begin{ejercicio}
        Se considera el problema de valores iniciales
        \[
            x' = \log(t - x^2), \quad x(1) = 0.
        \]
        ¿Existe una solución maximal definida en $\left]-\infty, +\infty\right[$?
    \end{ejercicio}

    \begin{ejercicio}
        El campo vectorial $X : \bb{R} \times \bb{R}^3 \to \bb{R}^3$, $(t, x) \mapsto X(t, x)$, es continuo y cumple
        \[
            \|X(t, x)\| \leq 3 \text{ si } t \in \bb{R}, \ \|x\| \leq 1.
        \]
        Prueba que si $x : \left]\alpha, \omega\right[ \to \bb{R}^3$ es una solución maximal de
        \[
            \dot{x} = X(t, x), \quad x(0) = 0
        \]
        entonces $\omega \geq \nicefrac{1}{3}$.
    \end{ejercicio}

    \newpage
    \setcounter{ejercicio}{0} % Reiniciar contador de ejercicios
    \noindent
    % \textbf{Solución.}

\end{document}